\documentclass{estilo}

\usepackage[spanish, es-nodecimaldot]{babel}
\usepackage{float}
\usepackage{amsmath}        % para los vectores columnas
\usepackage{amsfonts}       % para las negrita de pizarra
\usepackage{amssymb}        % para simbolos matematicos
\usepackage{mathtools}      % floor y ceil
\usepackage{multirow}       % para las tablas
\usepackage{dsfont}
\usepackage{listings}
\usepackage{xcolor}
\usepackage{enumitem,multicol,setspace}
\usepackage{ragged2e}       % para que no tire error con \justifying

\definecolor{codegreen}{rgb}{0,0.6,0}
\definecolor{codegray}{rgb}{0.5,0.5,0.5}
\definecolor{codepurple}{rgb}{0.58,0,0.82}
\definecolor{backcolour}{rgb}{0.95,0.95,0.92}
\definecolor{codeblue}{rgb}{0.16,0.36,0.67}
% config. para que el pseudocodigo tenga colorcitos.
\lstdefinestyle{mystyle}{
    backgroundcolor=\color{backcolour},   % color de fondo, ese lo ponemos?
    commentstyle=\color{codegreen},       % color de comentarios
    keywordstyle=\color{codepurple}\bfseries, % color de palabras clave (if, return, def)
    numberstyle=\tiny\color{codegray}, % color de números de línea
    stringstyle=\color{codeblue}, % color de strings
    basicstyle=\ttfamily\small, % tipografia monoespaciada
    breakatwhitespace=false,         
    breaklines=true,                 
    captionpos=b,                    
    keepspaces=true,                 
    numbers=right,
    numbersep=5pt,                  
    showspaces=false,                
    showstringspaces=false,
    showtabs=false,                  
    tabsize=4,
    language=Python % detecta automáticamente def, return, if, else, etc.
}
\lstset{style=mystyle}

\onehalfspacing % Interlineado 1.5
\setlength{\parskip}{1.2ex plus 0.4ex minus 0.2ex} % Espacio entre párrafos
\setlength{\parindent}{0pt} % Elimina la sangría obligatoria para que no quede raro 


\usepackage{hyperref} 
\hypersetup{
    hidelinks,
    pdfstartview={XYZ null null null},
    pdfview={XYZ null null null}
}

% bibliografía con BibLaTeX
\usepackage[style=authoryear-comp]{biblatex}
\addbibresource{referencias.bib}


\newcounter{subenum}[enumi] % para las multicolumnas
\renewcommand{\thesubenum}{\arabic{subenum}}
\usepackage[nomessages]{fp}
\FPeval\thecolwidth{round(1/4:4)}% Specify number of columns -> column width
\newcommand{\newitem}[1]{%
  \refstepcounter{subenum}%
  \parbox{\dimexpr\thecolwidth\linewidth-.5\columnsep}{%
    \makebox[\labelwidth][r]{(\thesubenum)\hspace*{\labelsep}}%
    #1}\hfill%
}

% notacion matemática complementaria
\usepackage{scalerel,stackengine} % para el sombrero
\stackMath
\newcommand\rhat[1]{%
\savestack{\tmpbox}{\stretchto{%
  \scaleto{%
    \scalerel*[\widthof{\ensuremath{#1}}]{\kern-.6pt\bigwedge\kern-.6pt}%
    {\rule[-\textheight/2]{1ex}{\textheight}}%WIDTH-LIMITED BIG WEDGE
  }{\textheight}% 
}{0.5ex}}%
\stackon[1pt]{#1}{\tmpbox}%
}
\parskip 1ex


\DeclarePairedDelimiter\ceil{\lceil}{\rceil}
\DeclarePairedDelimiter\floor{\lfloor}{\rfloor} 




\begin{document}
\nocite{*} % para la bibliografía.

\maketitle
\justifying

%\newpage
% el resumen va directo despues del titulo. no se usa caratula separada.
\vspace*{1.5cm}
\section{Resumen}

nashe

\vspace{0.4cm}
\begin{center}
    \vspace*{0.1cm} 
    \Huge\bfseries Índice
    \vspace{0.1cm}
\end{center}

\renewcommand{\contentsname}{} % elimina el indice chico que mete latex por defecto

% El comando automático ahora va a renderizar links cliqueables a tus archivos
\tableofcontents

\clearpage 


% aparece en el indice pero sin numero: \addcontentsline{toc}{section}{Conclusion}


\newpage
\section{El Problema}

nashe

\subsection{Introducción}

nashe nashe

\subsection{Planteamiento}

nashe nashe

\subsection{Consecuencias de la optimalidad esto lo pongo?}

nashe

\newpage
\section{Diseño del Algoritmo}

nashe

\subsection{Código}\label{rec:codigo} % para llamarlo después. en el analisis de la complejidad.

\begin{lstlisting}[mathescape=true, language=Python]
def acacodigo
def recons

\end{lstlisting}


\newpage
\section{Analizando el Algoritmo}

nashe


\newpage
\section{Mejoras}


nashe


\newpage
\section{Conclusión}

El desarrollo de este Trabajo Práctico nos permitió investigar y entender a fondo el poder de Programación Dinámica. Con el diseño de la ecuación de recurrencia y la matriz de memoization, se pudo resolver un escenario donde cada decisión presente desencadena un efecto dominó en el futuro, garantizando matemáticamente el máximo valor acumulado posible para Sophia en cada partida.

Al contrastar esta solución con el enfoque Greedy, mientras que solución avariciosa tenía un costo rapidísimo en $\mathcal{O}(n)$ usando memoria constante, falla al caer en ciertas ``trampas'' perdiendo la visión del problema completo frente a variabilidades grandes. En cambio, la Programación Dinámica exige pagar un precio alto: un sacrificio computacional cuadrático tanto en tiempo ($\mathcal{O}(n^2)$) como en espacio ($\mathcal{O}(n^2)$) para guardar la ``foto o recuerdos'' de la experiencia en una matriz. Pero nos da una garantía de optimalidad inquebrantable.

En conclusión, esta forma nos demuestra que la mejor solución global rara vez se encuentra eligiendo el camino más fácil y evidente en el presente. Resolver problemas complejos de verdad requieren la capacidad de sacrificar el beneficio inmediato, pensar en las ramificaciones del futuro, y tener la memoria suficiente para no recalcular los mismos errores del pasado.


\vfill %para que la bibliografia vaya para abajo.
\printbibliography[heading=bibintoc, title={Bibliografía}]

\end{document}
